\chapter{Mercados garantidos por terceiros}
\label{chapter:escrowed-market}

Uma garantia 2-de-3 coloca comprador, vendedor e moderador no mesmo endereço
multi-assinatura. Qualquer par pode autorizar um gasto; nenhum participante
pode fazê-lo sozinho. Esta propriedade permite liquidar uma compra por acordo
directo ou por uma decisão apoiada pelo moderador.

A garantia não é um contracto executado pelo consenso do Monero. O consenso
verifica uma transacção RingCT e a respectiva assinatura. Não conhece a ordem
de compra, a entrega, o prazo, as provas apresentadas numa disputa nem a
identidade do moderador. Esses elementos pertencem ao software do mercado e
ao protocolo de comunicação entre os participantes.

A análise histórica de integração com OpenBazaar identificou os problemas
operacionais que uma aplicação deste tipo teria de resolver
\cite{openbazaar-rbrunner-investigation}. O código Monero consultado implementa
as operações gerais de carteira multi-assinatura e o Sistema de Mensagens de
Multi-assinatura, MMS \cite{monero-source-v16,mms-manual,mms-project-proposal}.
Não implementa um mercado, um formato de ordem, um processo de arbitragem ou
uma política automática de reembolso.

\section{Estado da implementação}
\label{sec:escrow-implementation-status}

Na revisão de 14 de Agosto de 2026, a separação entre componentes
implementados e componentes externos é a seguinte:

\begin{center}
\small
\begin{tabular}{|p{4.2cm}|p{3.0cm}|p{7.1cm}|}
\hline
\textbf{Componente} & \textbf{Estado} & \textbf{Consequência} \\
\hline
Carteira \(M\)-de-\(N\), incluindo 2-de-3 &
Implementada, experimental &
Os três participantes podem criar um endereço comum e quaisquer dois podem
produzir uma transacção válida. \\
\hline
Exportação e importação de informação multi-assinatura &
Implementada, experimental &
As imagens de chave parciais e o material necessário à assinatura podem ser
sincronizados entre carteiras independentes. \\
\hline
Proposta, assinatura e submissão de transacções &
Implementada, experimental &
Uma proposta pode receber contribuições até atingir o limiar e ser difundida
como uma transacção RingCT comum. \\
\hline
MMS &
Implementado &
Organiza mensagens e estados da carteira; não decide disputas nem substitui a
autenticação dos participantes. \\
\hline
Ordem de compra, prova de entrega e decisão do moderador &
Não implementadas pelo Monero &
A aplicação tem de definir formatos, assinaturas, armazenamento e regras de
validação. \\
\hline
Reembolso unilateral após um prazo &
Não implementado &
Uma saída Monero não contém um ramo de script que altere o conjunto de
signatários depois de uma data. \\
\hline
Acesso do moderador apenas durante uma disputa &
Não implementado pela carteira multi-assinatura &
Os participantes partilham a chave privada de vista comum; o moderador pode
observar a carteira desde a configuração. \\
\hline
Multi-assinatura com FCMP++ &
Primitivas de referência, sem integração na carteira &
O ramo FCMP++ de \texttt{monero-oxide} contém algoritmos SAL com FROST; a
construção da carteira Monero termina em CLSAG distribuída, como indicado na secção
\ref{sec:multisig-fcmp-note}. \\
\hline
\end{tabular}
\end{center}

Consequentemente, este capítulo especifica uma utilização possível das
operações implementadas. Não afirma que exista uma aplicação de mercado
fornecida pelo projecto Monero.

\section{Modelo formal}
\label{sec:escrow-formal-model}

Sejam \(B\), \(V\) e \(M\) o comprador, o vendedor e o moderador. A política
de autorização é
\begin{align}
 \mathcal A
 =\bigl\{S\subseteq\{B,V,M\}:|S|\geq2\bigr\}.
 \label{eq:escrow-authorized-sets}
\end{align}
Para um conjunto de signatários \(S\), uma transacção pode ser concluída
quando
\begin{align}
 S\in\mathcal A
 \quad\Longleftrightarrow\quad
 |S|\geq 2.
 \label{eq:escrow-threshold}
\end{align}

O modelo fornece as propriedades seguintes, sob as hipóteses criptográficas
da carteira multi-assinatura:

\begin{enumerate}
 \item \textbf{Autorização:} um participante isolado não produz a assinatura
 necessária para gastar as saídas da carteira.

 \item \textbf{Liquidação por acordo:} \(B\) e \(V\) podem pagar ou reembolsar
 sem intervenção de \(M\).

 \item \textbf{Liquidação de uma disputa:} \(M\) pode cooperar com \(B\) ou
 com \(V\), segundo a decisão tomada fora da lista de blocos.

 \item \textbf{Exclusão entre resultados:} duas transacções alternativas que
 gastam as mesmas saídas publicam as mesmas imagens de chave; no máximo uma
 pode ser confirmada.

 \item \textbf{Ausência de disponibilidade garantida:} se menos de dois
 participantes fornecem material válido e actualizado, os fundos não podem
 ser gastos.
\end{enumerate}

A política não garante que o moderador seja imparcial, que a mercadoria exista
ou que as provas externas sejam verdadeiras. Também não impede uma coligação
de dois participantes. Por exemplo, \(\{V,M\}\) pode pagar ao vendedor sem o
consentimento do comprador, porque esse resultado é precisamente autorizado
por \eqref{eq:escrow-authorized-sets}.

\subsection{Fundos e resultados possíveis}

Se o valor confirmado na carteira de garantia é \(e\) e a taxa da
transacção de liquidação é \(f\), um resultado é o vector
\begin{align}
 \mathbf v=(v_B,v_V,v_M),
 \qquad
 v_B+v_V+v_M+f=e,
 \label{eq:escrow-value-allocation}
\end{align}
com \(v_B,v_V,v_M\geq0\). A componente \(v_M\) pode ser zero ou pode conter
uma taxa de moderação previamente acordada.

Três resultados comuns são:

\begin{align}
 \mathbf v_{\rm pagamento}
   &=(e-p-f,\ p,\ 0),\\
 \mathbf v_{\rm reembolso}
   &=(e-f,\ 0,\ 0),\\
 \mathbf v_{\rm div}
   &=(e-p'-m-f,\ p',\ m),
 \label{eq:escrow-common-outcomes}
\end{align}
onde \(p\) é o preço, \(p'\) é o montante atribuído ao vendedor numa decisão
parcial e \(m\) é a taxa paga ao moderador. Cada vector tem de ser verificado
pelos signatários antes da assinatura.

No nível RingCT, sejam \(C_j^{\rm in}\) os pseudo-compromissos das entradas e
\(C_k^{\rm out}\) os compromissos das saídas. A transacção aceite satisfaz
\begin{align}
 \sum_j C_j^{\rm in}
 =
 \sum_k C_k^{\rm out}+fH.
 \label{eq:escrow-ringct-balance}
\end{align}
A equação prova o equilíbrio dos compromissos; não prova que
\(\mathbf v\) corresponde ao contracto comercial.

\subsection{Identidade de uma ordem}

A aplicação deve fixar uma representação canónica da ordem. Um exemplo é
\begin{align}
 \mathsf{ord}={}&(
 \mathsf{version},\mathsf{id},B,V,M,
 \mathsf{description},p,m_{\max},\\
 &\mathsf{delivery\_terms},
 \mathsf{dispute\_deadline},
 \mathsf{evidence\_policy},
 \mathsf{refund\_address},
 \mathsf{payment\_address}).
 \label{eq:escrow-order-record}
\end{align}
A identidade criptográfica da ordem é
\begin{align}
 h_{\rm ord}
 =\mathcal H(
 \texttt{monero-escrow-order-v1}
 \mathbin\|\operatorname{enc}(\mathsf{ord})).
 \label{eq:escrow-order-hash}
\end{align}
A função \(\operatorname{enc}\) tem de ser determinística: a mesma ordem não
pode admitir duas sequências de bytes. Comprador, vendedor e moderador
confirmam \(h_{\rm ord}\), as chaves de autenticação e o endereço
multi-assinatura antes do financiamento.

As assinaturas da aplicação não devem reutilizar chaves ou valores aleatórios
do protocolo de gasto. Cada participante dispõe de uma chave de autenticação
separada e assina, pelo menos,
\begin{align}
 \sigma_i
 =\operatorname{Sign}_{a_i}(
 h_{\rm ord},
 K_{\rm escrow}^s,
 K_{\rm escrow}^v,
 \mathsf{role}_i).
 \label{eq:escrow-role-binding}
\end{align}
Esta assinatura liga o papel comercial à carteira concreta. A confirmação
independente das impressões digitais das chaves continua necessária.

\section{Configuração da carteira}
\label{sec:escrow-wallet-setup}

Cada ordem usa uma carteira 2-de-3 nova. Os três participantes executam o
protocolo interactivo da secção \ref{sec:m-of-n} e terminam com o mesmo par
público
\begin{align}
 (K_{\rm escrow}^s,K_{\rm escrow}^v).
\end{align}
A construção requer as rondas de troca de chaves da implementação. Publicar
uma chave-base de produto ou pré-seleccionar um moderador não permite ao
comprador calcular unilateralmente o endereço multi-assinatura actual.

Cada participante deve:

\begin{enumerate}
 \item criar uma carteira nova, sem histórico e sem reutilizar uma cópia de
 outra carteira convertida;

 \item autenticar as mensagens de troca de chaves;

 \item concluir todas as rondas exigidas pela carteira;

 \item confirmar por um canal independente o endereço final e
 \(h_{\rm ord}\);

 \item conservar cópias seguras da sua própria carteira e do histórico de
 mensagens necessário à recuperação.
\end{enumerate}

As três carteiras conhecem a chave privada de vista comum. Portanto, todos os
participantes podem detectar as saídas recebidas e descodificar os montantes.
Depois de trocarem as imagens de chave parciais, podem também acompanhar quais
dessas saídas foram gastas. A confidencialidade dos montantes perante
observadores externos não implica confidencialidade entre \(B\), \(V\) e \(M\).

\section{Financiamento}
\label{sec:escrow-funding}

O comprador financia o endereço apenas depois da confirmação da configuração.
A mensagem de financiamento fixa
\begin{align}
 \mathsf{fund}=
 (h_{\rm ord},\mathsf{txid}_{\rm fund},e,n_{\rm conf}),
 \label{eq:escrow-funding-record}
\end{align}
onde \(n_{\rm conf}\) é a política de confirmações acordada. O vendedor não
deve tratar uma transacção apenas observada na memória temporária como
liquidação definitiva.

Depois das confirmações, cada carteira sincroniza a saída. Os participantes
trocam informação multi-assinatura para construir as imagens de chave
completas. Para uma saída com chave \(P\) e chave privada de utilização única
distribuída
\begin{align*}
 x=d+\sum_{j\in\mathcal I}x_j,
\end{align*}
as parcelas combinam-se em
\begin{align}
 I=x\mathcal H_p(P)
  =d\mathcal H_p(P)+
    \sum_{j\in\mathcal I}x_j\mathcal H_p(P).
 \label{eq:escrow-key-image}
\end{align}
Sem a importação suficiente, uma carteira pode não conhecer correctamente o
estado gasto da saída e pode não dispor do material necessário à proposta.

O montante \(e\) deve incluir a taxa prevista para a liquidação e uma margem
definida pela aplicação. Uma estimativa insuficiente não pode ser corrigida
por uma saída inexistente; pode exigir uma entrada adicional e, por
consequência, nova sincronização e nova proposta.

\section{Registo autenticado das mensagens}
\label{sec:escrow-authenticated-log}

A decisão do moderador depende de dados que não estão na lista de blocos. A
aplicação deve produzir um registo verificável. Seja a mensagem \(j\)
\begin{align}
 m_j=(
 h_{\rm ord},j,\mathsf{type}_j,\mathsf{payload}_j,h_{j-1}),
 \qquad
 h_j=\mathcal H(\operatorname{enc}(m_j)),
 \label{eq:escrow-message-chain}
\end{align}
e seja
\begin{align}
 \sigma_j=\operatorname{Sign}_{a_{\mathsf{sender}(j)}}(h_j).
 \label{eq:escrow-message-signature}
\end{align}
O número de sequência e a hash anterior detectam omissões, reordenações e
substituições dentro da transcrição apresentada. Não demonstram que um facto
físico, como uma entrega, aconteceu; demonstram quem assinou determinada
afirmação e qual era a sua posição no registo.

Os campos \(\mathsf{type}_j\) devem distinguir pelo menos:

\begin{itemize}
 \item aceitação ou rejeição da ordem;
 \item confirmação do financiamento;
 \item declaração de expedição ou prestação;
 \item confirmação de recepção;
 \item pedido de pagamento ou de reembolso;
 \item abertura de disputa;
 \item apresentação de prova;
 \item decisão do moderador;
 \item referência à transacção final.
\end{itemize}

A aplicação deve definir limites de tamanho, codificação, relógio, política de
retenção e tratamento de mensagens duplicadas. A encriptação do canal protege
o conteúdo em trânsito; as assinaturas autenticam o conteúdo depois de
descodificado. São propriedades diferentes.

\section{Liquidação sem disputa}
\label{sec:escrow-cooperative-settlement}

Depois de o comprador aceitar o resultado, \(B\) e \(V\) escolhem um vector
\(\mathbf v\) de \eqref{eq:escrow-common-outcomes}. Um deles constrói uma
proposta que fixa:

\begin{align}
 \mathsf{proposal}=(
 h_{\rm ord},
 \mathsf{txid}_{\rm fund},
 \mathbf v,
 f,
 \mathsf{destinations},
 \mathsf{change},
 \mathsf{proposal\_id}).
 \label{eq:escrow-proposal}
\end{align}

Antes de assinar, cada participante verifica directamente na proposta:

\begin{enumerate}
 \item que as entradas pertencem à carteira da ordem;
 \item que os destinos correspondem aos endereços autenticados;
 \item que os montantes realizam \(\mathbf v\);
 \item que a taxa é aceitável;
 \item que não existem saídas adicionais;
 \item que a proposta ainda não foi substituída ou cancelada.
\end{enumerate}

A sequência de comandos é a da secção \ref{sec:multisig-cli}: sincronizar
informação multi-assinatura, criar uma proposta, adicionar duas contribuições
e submeter a transacção completa. O ficheiro da proposta é parte da
transcrição criptográfica de gasto, não uma descrição comercial da ordem.

Os dados exportados e os valores aleatórios de assinatura são de uso único.
Uma nova proposta, uma alteração de taxa ou uma nova exportação exige que os
participantes abandonem de forma explícita o material anterior. Duas propostas
concorrentes não devem partilhar valores aleatórios privados.

\section{Disputa e decisão}
\label{sec:escrow-dispute}

Uma disputa começa com uma mensagem assinada que contém \(h_{\rm ord}\), o
resultado pedido e a hash do registo apresentado. O moderador verifica:

\begin{enumerate}
 \item as assinaturas e a continuidade do registo;
 \item a correspondência entre a ordem e a carteira financiada;
 \item a política de prova definida na ordem;
 \item as afirmações de ambas as partes;
 \item o saldo e o estado das saídas da carteira;
 \item os endereços que receberão cada componente do resultado.
\end{enumerate}

A decisão tem uma forma canónica, por exemplo
\begin{align}
 \mathsf{decision}=
 (h_{\rm ord},h_{\rm log},\mathbf v,\mathsf{reason},
 \mathsf{decision\_id}),
 \label{eq:escrow-decision}
\end{align}
e o moderador publica
\begin{align}
 \sigma_M^{\rm decision}
 =\operatorname{Sign}_{a_M}(
 \mathcal H(\operatorname{enc}(\mathsf{decision}))).
 \label{eq:escrow-decision-signature}
\end{align}

A decisão não move fundos. Para a executar, o moderador e a parte que a aceita
sincronizam as carteiras, verificam uma proposta que realiza
\(\mathbf v\) e fornecem as duas contribuições necessárias. Os conjuntos
\(\{B,M\}\) e \(\{V,M\}\) têm exactamente o mesmo poder criptográfico; a
diferença entre reembolso e pagamento está nos destinos e montantes da
proposta.

Se o comprador e o vendedor resolverem a disputa antes da liquidação, podem
assinar directamente uma proposta substituta. A aplicação deve marcar a
decisão anterior como não executada e registar a transacção efectivamente
confirmada.

\subsection{Resultados concorrentes}

Considere duas transacções válidas \(T_1\) e \(T_2\) que gastam pelo menos uma
mesma saída \(P_j\). A imagem de chave dessa entrada é determinística:
\begin{align}
 I_j=x_j\mathcal H_p(P_j).
 \label{eq:escrow-double-spend-image}
\end{align}
A política de consenso rejeita a segunda ocorrência confirmada de \(I_j\).
Assim,
\begin{align}
 T_1\text{ confirmada}
 \quad\Longrightarrow\quad
 T_2\text{ inválida relativamente ao estado resultante}.
 \label{eq:escrow-outcome-exclusion}
\end{align}

Antes da confirmação, propostas incompatíveis podem circular
simultaneamente. A aplicação não deve tratar uma assinatura parcial, uma
transacção na memória temporária ou uma decisão do moderador como confirmação
final. A referência final é a transacção aceite na lista de blocos com a
profundidade exigida pela política da ordem.

\section{Prazos e disponibilidade}
\label{sec:escrow-timeouts}

Um prazo comercial não altera \(\mathcal A\). Depois da data indicada na
ordem, o comprador continua sem poder reembolsar sozinho e o vendedor
continua sem poder pagar-se sozinho. O campo de bloqueio temporal de uma
transacção Monero não implementa uma condição do tipo
\begin{align*}
 \text{antes de }t:\ 2\text{-de-3},
 \qquad
 \text{depois de }t:\ B\text{ sozinho}.
\end{align*}

Uma transacção de reembolso completamente assinada com antecedência é apenas
um gasto alternativo já autorizado por dois participantes. Não constitui um
ramo de consenso activado no futuro e pode competir com outras transacções
que usem as mesmas entradas. Alterações de taxa, reorganizações, entradas
adicionais ou perda do ficheiro também podem impedir o comportamento esperado.

A disponibilidade pode ser resumida por
\begin{align}
 \operatorname{live}(S)
 =
 \begin{cases}
 1,&|S\cap\{B,V,M\}|\geq2
      \text{ e os dados estão sincronizados},\\
 0,&\text{caso contrário}.
 \end{cases}
 \label{eq:escrow-liveness}
\end{align}
O moderador é, portanto, parte da disponibilidade durante uma disputa. A
aplicação deve definir substituição de moderadores antes do financiamento ou
aceitar que a indisponibilidade permanente de \(M\), combinada com desacordo
entre \(B\) e \(V\), bloqueia os fundos.

\section{Privacidade}
\label{sec:escrow-privacy}

Para um observador da lista de blocos, o financiamento e a liquidação têm a
forma normal de transacções Monero. A política 2-de-3, os papéis dos
participantes e o conteúdo da disputa não são campos públicos de consenso.

Entre os participantes, a privacidade é menor:

\begin{itemize}
 \item todos conhecem o endereço multi-assinatura e a chave privada de vista
 comum;
 \item todos podem associar o financiamento a \(h_{\rm ord}\);
 \item os signatários de uma proposta conhecem entradas, destinos, montantes
 e taxa;
 \item o moderador recebe o registo e as provas definidos pela política da
 aplicação;
 \item metadados de rede e de comunicação podem ligar identidades mesmo
 quando a lista de blocos não o faz.
\end{itemize}

A proposta histórica de entregar ao moderador a capacidade de vista apenas
quando surge uma disputa não corresponde à carteira multi-assinatura
implementada. Essa propriedade exigiria uma construção de chaves ou uma
camada de custódia de dados diferente. Encriptar o registo comercial até à
disputa pode limitar o acesso às mensagens, mas não retira ao moderador a
chave de vista da carteira.

Cada ordem deve usar novas chaves e novos endereços de destino. A reutilização
de uma carteira 2-de-3 entre ordens permite aos três participantes relacionar
os respectivos financiamentos e liquidações.

\section{Relação com FCMP++}
\label{sec:escrow-fcmp}

A política comercial \(\mathcal A\) não depende de CLSAG. Se a carteira
multi-assinatura for adaptada a FCMP++, a configuração continua a distribuir
o segredo de gasto e a liquidação continua a exigir duas contribuições. A
diferença está na prova produzida para cada entrada:

\begin{align*}
 \text{CLSAG distribuída}
 \quad\longrightarrow\quad
 \text{pertença FCMP}+\text{SAL distribuída}.
\end{align*}

A lista de blocos continuaria sem publicar \(B,V,M\) ou o limiar. Contudo, a
implementação precisa de construir conjuntamente o marcador de ligação, o
produto \(z=xr_i\), os compromissos e as respostas SAL sem reconstruir a chave
privada completa. O ramo de referência contém primitivas distribuídas para
essa prova, mas a sua integração no fluxo da carteira Monero não existe na
revisão consultada. A secção \ref{sec:multisig-fcmp-note} distingue esses dois
níveis.

\section{Requisitos mínimos de uma aplicação}
\label{sec:escrow-application-requirements}

Uma aplicação que pretenda usar esta construção deve implementar e testar,
pelo menos:

\begin{enumerate}
 \item codificação canónica e assinatura de ordens, decisões e mensagens;
 \item autenticação dos três participantes e confirmação independente do
 endereço;
 \item uma carteira nova por ordem;
 \item armazenamento cifrado e cópias de segurança de cada carteira;
 \item sincronização de informação multi-assinatura antes de cada proposta;
 \item validação local de entradas, destinos, montantes, troco e taxa;
 \item controlo de uso único dos valores aleatórios e dos ficheiros de
 proposta;
 \item uma política explícita de confirmações e reorganizações;
 \item um formato verificável de prova e decisão;
 \item tratamento de indisponibilidade e de fundos bloqueados;
 \item protecção dos canais de comunicação e dos metadados;
 \item testes contra propostas concorrentes, mensagens repetidas, estados
 divergentes e participantes maliciosos.
\end{enumerate}

Estas funções não são acessórios de interface. São parte do modelo de
segurança do mercado. A assinatura 2-de-3 limita quem pode autorizar a
transacção; a aplicação determina qual transacção cada participante deve
autorizar.

\section{Resumo}

Uma garantia Monero 2-de-3 usa a política
\[
 \mathcal A
 =\{\{B,V\},\{B,M\},\{V,M\},\{B,V,M\}\}.
\]
A carteira multi-assinatura implementa o limiar e produz uma transacção RingCT
comum. O mercado tem de implementar a ordem, a autenticação, o registo de
provas, a decisão, os prazos e a verificação da proposta.

Qualquer par pode liquidar; um participante isolado não pode. Duas
liquidações incompatíveis que gastem as mesmas saídas têm imagens de chave
iguais, pelo que no máximo uma pode ser confirmada. Não existe reembolso
unilateral por prazo, acesso condicional nativo do moderador nem fluxo de
mercado fornecido pelo Monero. A segurança operacional depende de carteiras
independentes, mensagens autenticadas, sincronização correcta e valores
aleatórios de uso único.
